\subsection{Do counterfeit samples from different models or problems differ significantly?} \label{sec:analyzing-counterfeit}
\subsubsection{Is it easier for models to understand counterfeit samples from problems it finds easier?}
Intuitively, if a given programming problem is easy for a model to solve (meaning it has a high pass@1 rate), we might believe models understand how to solve that problem. If that is the case, then models should be able to better understand both correct and counterfeit samples for that problem. To test if this was true, we bucketed problem difficulties into easy, medium, and hard by a model's pass@1 on that problem. We then calculated the average correctness checking, verification, and repair accuracy for each of the tasks for each problem difficulty bucket. Fig. \ref{fig:difficulty-sample} shows a representative subset of these results on HumanEval.

\begin{figure}[ht!]
    \centering
    \begin{subfigure}[b]{0.31\textwidth}
        \centering
        \includegraphics[width=\textwidth]{figures_v2/main/difficulty_barplots/humaneval_correctness_sample.pdf}
        \subcaption{Correctness Checking}
    \end{subfigure}
    \hfill
    \begin{subfigure}[b]{0.31\textwidth}
        \centering
        \includegraphics[width=\textwidth]{figures_v2/main/difficulty_barplots/humaneval_execution_sample.pdf}
        \subcaption{Execution}
    \end{subfigure}
    \hfill
    \begin{subfigure}[b]{0.31\textwidth}
        \centering
        \includegraphics[width=\textwidth]{figures_v2/main/repair/humaneval/repair_4x4.pdf}
        \subcaption{Repair}
    \end{subfigure}
    \caption{Mean accuracy of correctness checking, execution, and repair on HumanEval (error bars are 1 std). A problem's difficulty is determined by the pass@1 of the evaluation model. In general, correctness checking accuracy is uncorrelated with problem difficulty, while both execution and repair accuracies are only very weakly correlated.}
    \label{fig:difficulty-sample}
\end{figure}

Surprisingly, we found that 1) correctness checking accuracies were relatively uncorrelated with problem difficulty, while 2) execution ability and the success rate of repair exhibited a modest amount of correlation with problem difficulty.
We found that these trends were generally robust across HumanEval, LeetCode, and ODEX, although the amount of correlation exhibited in the repair task varied (but was, at best, modest).

To get a more precise measurement of these relationships, we calculated the Pearson correlation between generation performance and each of correctness checking, execution prediction, and repair performance across problems. The resulting histogram showing the distribution of correlations for each task is shown in Fig. \ref{fig:correlation-with-generation}.

\begin{figure}[H]
\centering
\includegraphics[width=0.75\columnwidth]{figures_v2/main/correlation/correlation_with_generation.pdf}
\caption{Correlation between generation and each of our three tasks}
\label{fig:correlation-with-generation}
\end{figure}

\subsubsection{Do models perceive their own samples differently?}
For a given model, its counterfeit samples had a high enough log-likelihood to be generated by the model, so one may hypothesize that models might have a harder time than other models at distinguishing their own counterfeit samples. In Fig. \ref{fig:self-understanding-heatmaps}, we plotted heatmaps showing the performance of various models on datasets generated by other models for the correctness checking task (left) and execution prediction task for HumanEval (right). For both tasks, the relative performance of different models was similar across datasets, and we found no evidence that models faltered more on their own samples. This suggests that counterfeit samples may be general: those from one model are generally difficult for other models to understand as well.
% It would have been really nice to do the perplexity experiment :(

\begin{figure}[H]
    \centering
    \hspace*{\fill}
    \begin{subfigure}[b]{0.4\columnwidth}
        \centering
        \includegraphics[width=\textwidth]{figures_v2/main/accuracy_heatmaps/odex_verification.pdf}
        \subcaption{Correctness (ODEX)}
    \end{subfigure}
    \hspace*{\fill}
    \begin{subfigure}[b]{0.4\columnwidth}
        \centering
        \includegraphics[width=\textwidth]{figures_v2/main/accuracy_heatmaps/humaneval_execution.pdf}
        \subcaption{Execution (HumanEval)}
    \end{subfigure}
    \hspace*{\fill}
    \caption{Heatmap of accuracies for correctness checking and execution prediction. Models do not seem to perceive their own generations differently from those of other models.}
    \label{fig:self-understanding-heatmaps}
\end{figure}

\subsubsection{Do stronger models generate harder counterfeit samples?}
One might also expect that counterfeit samples of stronger models are harder to verify than those of weaker models, as stronger models are less likely to generate obvious mistakes. In Fig. \ref{fig:average-counterfeit-accuracy}, we compared the average scores of two tasks on counterfeit samples for datasets generated by stronger (DS-33B, CL-34B) and weaker (CL-7B, SC) models. Note that this is the same as Fig. \ref{fig:self-understanding-heatmaps} with each row aggregated and filtered to only include counterfeit samples. Since there did not seem to be a significant difference between the difficulties, models of all strengths could be used to generate counterfeit samples that were challenging for models to understand.

\begin{figure}[ht!]
    \centering
    \begin{subfigure}[b]{0.49\columnwidth}
        \centering
        \includegraphics[width=\textwidth]{figures_v2/main/average_model_accuracy/odex_verification.pdf}
        \subcaption{Correctness (ODEX)}
    \end{subfigure}
    \hfill
    \begin{subfigure}[b]{0.49\columnwidth}
        \centering
        \includegraphics[width=\textwidth]{figures_v2/main/average_model_accuracy/humaneval_execution.pdf}
        \subcaption{Execution (HumanEval)}
    \end{subfigure}
    \caption{Accuracy of correctness and execution tasks on counterfeit samples. Both strong and weak models are able to generate difficult counterfeit samples.}
    \label{fig:average-counterfeit-accuracy}
\end{figure}
